\subsubsection{Weaknesses of Variational Algorithms}\label{sec:weaknesses-of-variational-algorithms}

As discussed in the previous section, VQAs seem to be the perfect solution for NISQ devices. However, \textbf{VQAs are not perfect}. We use them \textbf{because they are currently the best option for NISQ hardware}, not because they are theoretically superior to all quantum algorithms.

\highspace
The main weaknesses of VQAs are:
\begin{enumerate}
    \item They are \textbf{approximate methods}.
    \item The iterative process is \textbf{slow and resource-intensive}.
    \item They offer \textbf{few formal guarantees}.
    \item The \textbf{classical optimizer may fail} to find good parameters.
\end{enumerate}
Let's understand each of them.

\highspace
\begin{flushleft}
    \textcolor{Red2}{\faIcon{bullseye} \textbf{They Are Approximate Methods}}
\end{flushleft}
Unlike algorithms such as Shor's algorithm (we will discuss later), VQAs generally do \textbf{not} guarantee the exact solution. Instead, they \hl{search for the \textbf{best solution found}} rather than the \textbf{true optimal solution}. For example, in VQE, the algorithm might return an energy value that is very close to the ground state energy but not exact. Similarly, in QAOA, the algorithm may find a solution that is close to the optimal Max-Cut but not necessarily the best possible one. This is why QAOA is called the \textbf{Quantum Approximate Optimization Algorithm}. We will discuss QAOA in more detail in the next section.

\highspace
\begin{flushleft}
    \textcolor{Red2}{\faIcon{hourglass-half} \textbf{The Iterative Process is Slow and Resource-Intensive}}
\end{flushleft}
This is perhaps the most practical drawback. The optimization loop (\autopageref{sec:iterative-optimization-loop}) requires many iterations, and each iteration involves preparing the quantum state, executing the circuit, measuring, and often repeating the circuit many times to estimate expectation values.

\highspace
This makes \hl{VQAs computationally expensive and unlikely to show a quantum advantage on current devices. The total number of circuit executions can become enormous, making the process slow and resource-intensive.} But \emph{\textbf{why is it resource-intensive?}} Because one \hl{evaluation of the cost function is usually \textbf{not} one execution.} For example, in VQE, estimating the expectation value of the Hamiltonian requires many repeated measurements. So one optimization step actually involves multiple circuit executions, and when multiplied by hundreds or thousands of optimization iterations, the total number of circuit executions can become enormous.

\newpage

\begin{flushleft}
    \textcolor{Red2}{\faIcon{question-circle} \textbf{Few Formal Guarantees}}
\end{flushleft}
Classical optimization is difficult. Even if the \textbf{ansatz is expressive} and the \textbf{cost function is correct}, there is \hl{no guarantee that the optimizer will reach the global optimum.} It may stop too early or converge to a poor solution. Unlike algorithms such as Shor's algorithm, which are mathematically guaranteed to produce the correct answer (assuming ideal hardware), VQAs often lack such guarantees.

\highspace
\begin{flushleft}
    \textcolor{Red2}{\faIcon{exclamation-triangle} \textbf{The Optimizer May Fail}}
\end{flushleft}
This is closely related to the previous point. If the cost landscape has local minima, the \hl{optimizer may converge to a \textbf{local minimum} instead of the global minimum}. Since the gradient becomes zero at a local minimum, the optimizer may incorrectly conclude that it has found the best solution, even though a better solution exists elsewhere. This is why optimizer choice is important and why optimization itself is an active research area. However, this weakness is not a fundamental limitation of VQAs, but rather a classical optimization problem. In fact, the next section will discuss a specific reason why the classical optimizer may fail: \textbf{barren plateaus}.

\highspace
In general, all these weaknesses are consequences of one fundamental fact: \textbf{VQAs are optimization problems}. \hl{Finding the best parameters is often the hardest part of the algorithm.} Despite these drawbacks, VQAs remain the \textbf{most practical algorithms for current NISQ devices} because they can be executed on noisy, shallow quantum hardware.